\documentclass[12pt]{book} % Change to book
\usepackage{amsmath} % for mathematical expressions
\usepackage{xcolor} % for colored text
\usepackage{graphicx} % for including images
\usepackage{url}
\usepackage{booktabs} % for better-looking tables
\usepackage{makecell}
\usepackage{indentfirst}
\usepackage{algorithm}
\usepackage{algpseudocode}
\usepackage{amssymb}
\usepackage{placeins} % For \FloatBarrier
\usepackage{listings}
% Define the language style for Python code
\lstset{
    language=Python,
    basicstyle=\ttfamily,
    keywordstyle=\color{blue},
    stringstyle=\color{red},
    commentstyle=\color{gray},
    morecomment=[l][\color{magenta}]{\#},
    frame=single,
    breaklines=true
}
\begin{document}
\chapter{Scheduling: First Things First}

\section{Introduction to Scheduling}

Scheduling is a fundamental concept in computer science and operations research, concerned with efficiently allocating resources over time to accomplish tasks or jobs. It plays a crucial role in various domains, including operating systems, manufacturing, project management, and networking. In this section, we will delve into the definition, importance, and applications of scheduling techniques, followed by an overview of common scheduling problems.

\subsection{Definition and Importance}

Scheduling involves the allocation of resources such as processors, machines, or personnel to tasks or jobs over time, with the objective of optimizing certain criteria such as throughput, latency, response time, or resource utilization. Mathematically, scheduling can be represented as an optimization problem, where the goal is to find an assignment of resources to tasks that satisfies given constraints and optimizes an objective function.

One of the fundamental scheduling problems is the \textbf{single-machine scheduling problem}, where a set of jobs with given processing times and deadlines must be scheduled on a single machine to minimize a certain criterion, such as the total completion time or maximum lateness. Mathematically, the single-machine scheduling problem can be formulated as follows:

Given \(n\) jobs with processing times \(p_1, p_2, \ldots, p_n\) and deadlines \(d_1, d_2, \ldots, d_n\), schedule the jobs on a single machine to minimize an objective function \(f\) subject to certain constraints.

The importance of scheduling techniques lies in their ability to improve resource utilization, minimize response times, meet deadlines, and enhance overall system efficiency. Efficient scheduling algorithms are crucial for optimizing the performance of various systems and processes, leading to cost savings, improved productivity, and enhanced customer satisfaction.

\subsection{Applications in Computing and Beyond}

Scheduling techniques find applications in a wide range of computing and non-computing domains:

\begin{itemize}
    \item \textbf{Operating Systems:} Scheduling algorithms are used to manage the execution of processes or threads on a computer system's CPU. Examples include round-robin scheduling, shortest job first, and priority scheduling.
    \item \textbf{Manufacturing:} In manufacturing systems, scheduling is used to optimize production processes, allocate resources such as machines and workers, and minimize production costs. Examples include job shop scheduling and flow shop scheduling.
    \item \textbf{Networking:} Scheduling algorithms are employed in network protocols to manage the transmission of data packets, allocate network resources, and ensure fair access to network bandwidth. Examples include packet scheduling in routers and traffic scheduling in wireless networks.
    \item \textbf{Project Management:} Scheduling techniques are utilized in project management software to plan and allocate resources, schedule tasks, and track project progress. Examples include critical path method (CPM) and program evaluation and review technique (PERT).
\end{itemize}

\subsection{Overview of Scheduling Problems}

Scheduling problems can be classified based on various criteria, including the number of resources, the nature of the tasks, and the optimization criteria. Common scheduling problems include:

\begin{itemize}
    \item \textbf{Single-Machine Scheduling:} In this problem, a set of jobs must be processed on a single machine subject to certain constraints such as processing times and deadlines.
    \item \textbf{Parallel-Machine Scheduling:} This problem involves scheduling jobs on multiple identical or unrelated machines to optimize criteria such as makespan or total completion time.
    \item \textbf{Flow Shop Scheduling:} In flow shop scheduling, jobs must pass through a sequence of machines, each with its own processing time, to minimize the total completion time or makespan.
    \item \textbf{Job Shop Scheduling:} Job shop scheduling deals with scheduling jobs that require processing on multiple machines, each with different processing times and sequence-dependent setup times.
    \item \textbf{Resource-Constrained Project Scheduling:} This problem arises in project management, where tasks with precedence relations must be scheduled subject to resource constraints and project deadlines.
\end{itemize}

These scheduling problems are often NP-hard, requiring the development of efficient heuristic or approximation algorithms to find near-optimal solutions in a reasonable amount of time.


\section{Theoretical Foundations of Scheduling}

Scheduling is a fundamental problem in computer science and operations research, with applications in various domains such as manufacturing, transportation, and computing systems. This section provides an overview of the theoretical foundations of scheduling, including the classification of scheduling problems, optimality criteria, complexity measures, and computability.

\subsection{Classification of Scheduling Problems}

Scheduling problems can be classified based on various criteria, including the characteristics of the jobs or tasks to be scheduled, the processing environment, and the scheduling objectives. This classification helps in understanding the nature of different scheduling problems and developing appropriate solution approaches.

\textbf{Characteristics of Jobs or Tasks}

One common classification criterion is based on the characteristics of the jobs or tasks to be scheduled. These characteristics may include the following:

\begin{itemize}
    \item \textbf{Processing times:} Some scheduling problems involve jobs with deterministic processing times, where the processing time of each job is known in advance. In contrast, other problems may involve jobs with stochastic processing times, where the processing time of each job follows a probability distribution.
    
    \item \textbf{Release times and due dates:} In many scheduling problems, jobs have release times indicating when they become available for processing, and due dates indicating when they must be completed. Problems that consider release times and due dates are often referred to as "dynamic" scheduling problems, as they involve scheduling decisions over time.
    
    \item \textbf{Precedence constraints:} Some scheduling problems involve tasks that have precedence constraints, meaning that certain tasks must be completed before others can start. These problems are often referred to as "job shop" or "flow shop" scheduling problems and are common in manufacturing and production environments.
\end{itemize}

\textbf{Processing Environment}

Another classification criterion is based on the processing environment in which the jobs are scheduled. Common processing environments include the following:

\begin{itemize}
    \item \textbf{Single-machine scheduling:} In single-machine scheduling problems, there is only one machine available for processing the jobs. The objective is to determine the order in which the jobs should be processed on the single machine to optimize a certain criterion, such as makespan or total completion time.
    
    \item \textbf{Parallel machine scheduling:} In parallel machine scheduling problems, there are multiple identical machines available for processing the jobs. The objective is to assign the jobs to the machines in a way that minimizes a certain criterion, such as total completion time or maximum lateness.
    
    \item \textbf{Open-shop scheduling:} In open-shop scheduling problems, there are multiple machines available, and each job must be processed on a specific sequence of machines. The objective is to determine the processing sequence for each job to optimize a certain criterion, such as makespan or total completion time.
\end{itemize}

\textbf{Scheduling Objectives}

Scheduling problems can also be classified based on the scheduling objectives or criteria that are optimized. Common scheduling objectives include the following:

\begin{itemize}
    \item \textbf{Minimization of makespan:} The makespan is the total time required to complete all jobs in a schedule. Minimizing the makespan is a common objective in scheduling problems, as it ensures efficient resource utilization and timely job completion.
    
    \item \textbf{Minimization of total completion time:} The total completion time is the sum of the completion times of all jobs in a schedule. Minimizing the total completion time ensures that jobs are completed as early as possible, leading to shorter lead times and improved customer satisfaction.
    
    \item \textbf{Minimization of maximum lateness:} The maximum lateness is the maximum amount of time by which a job's completion time exceeds its due date. Minimizing the maximum lateness is important in scheduling problems with strict deadlines, as it ensures that jobs are completed on time to avoid penalties or disruptions.
\end{itemize}

By considering these classification criteria, scheduling problems can be analyzed and solved more effectively, leading to improved efficiency and performance in various application domains.

\subsection{Optimality Criteria}

Optimality criteria in scheduling problems determine the desired outcome or objective to be optimized. Common criteria include minimizing makespan, total completion time, maximum lateness, or total weighted completion time. 

Let $n$ be the number of jobs and $m$ be the number of machines. The scheduling problem can be formulated as follows:

\[
\text{Minimize } f(x_1, x_2, ..., x_n)
\]
subject to:
\[
g(x_1, x_2, ..., x_n) \leq b
\]
where $x_i$ represents the assignment of job $i$ to a machine, and $f$ is the objective function to be minimized.

\subsection{Complexity and Computability}

The complexity of scheduling problems is often analyzed using computational complexity theory. Let's consider a specific example: the scheduling problem with \(n\) jobs and \(m\) machines, where each job \(j\) has a processing time \(p_j\) and needs to be processed on one of the machines without preemption.

One common complexity measure is the number of operations required to find an optimal schedule. In this case, if we denote \(O(n, m)\) as the number of operations required to find an optimal schedule for \(n\) jobs and \(m\) machines, we can formulate the following proposition:

\textbf{Proposition 1:} The number of operations required to find an optimal schedule for the scheduling problem with \(n\) jobs and \(m\) machines is \(O(n \cdot m)\).

\textbf{Proof:} 
To find an optimal schedule, we need to consider all possible assignments of jobs to machines. Since each job can be assigned to any of the \(m\) machines, there are \(m^n\) possible combinations of job assignments. Therefore, the number of operations required to find an optimal schedule is \(O(n \cdot m)\).

The computability of scheduling problems refers to the feasibility of finding an optimal solution within a reasonable amount of time. Many scheduling problems are NP-hard, meaning that no polynomial-time algorithm exists unless P=NP. 

Let's consider the classical problem of flow shop scheduling, where \(n\) jobs need to be processed on \(m\) machines in a fixed sequence to minimize the makespan. This problem is known to be NP-hard.

\textbf{Theorem 1:} The flow shop scheduling problem is NP-hard.

\textbf{Proof:} 
We can reduce the problem of flow shop scheduling to the problem of traveling salesman, which is a well-known NP-hard problem. Given an instance of the traveling salesman problem with a set of cities and distances between them, we can construct an equivalent instance of the flow shop scheduling problem where each city corresponds to a job and each distance corresponds to a processing time on a machine. Therefore, since the traveling salesman problem is NP-hard, the flow shop scheduling problem is also NP-hard.


\section{Scheduling Algorithms}

Scheduling algorithms are essential in operating systems and computer science, determining the order in which tasks or processes are executed on a computing system. Different scheduling algorithms prioritize tasks differently based on various criteria such as waiting time, execution time, priority, etc. In this section, we will explore several common scheduling algorithms, including First-Come, First-Served (FCFS), Shortest Job First (SJF), Priority Scheduling, and Round-Robin Scheduling.

\subsection{First-Come, First-Served (FCFS)}

First-Come, First-Served (FCFS) is one of the simplest scheduling algorithms, where tasks are executed in the order they arrive in the system. It follows a non-preemptive approach, meaning once a process starts executing, it continues until completion without interruption.

Let's denote $n$ as the total number of processes, $AT_i$ as the arrival time of process $i$, $BT_i$ as the burst time of process $i$, and $CT_i$ as the completion time of process $i$. The completion time for each process can be calculated using the formula:

\[
CT_i = \sum_{j=1}^{i} BT_j
\]

The average waiting time ($\text{AWT}$) for all processes can be computed as the total waiting time divided by the number of processes ($n$):

\[
\text{AWT} = \frac{\sum_{i=1}^{n} WT_i}{n}
\]

where $WT_i$ is the waiting time for process $i$, calculated as:

\[
WT_i = CT_i - AT_i - BT_i
\]

Now, let's provide the algorithm for the FCFS technique and explain it in detail:
\FloatBarrier
\begin{algorithm}[H]
\caption{First-Come, First-Served (FCFS) Scheduling}
\begin{algorithmic}[1]
\State Sort processes based on arrival time $AT_i$
\State $CT_1 \gets AT_1 + BT_1$
\For{$i \gets 2$ to $n$}
    \State $CT_i \gets CT_{i-1} + BT_i$
\EndFor
\end{algorithmic}
\end{algorithm}
\FloatBarrier
The FCFS algorithm simply iterates through the processes in the order of their arrival and calculates the completion time for each process sequentially.

\subsection{Shortest Job First (SJF)}

Shortest Job First (SJF) is a scheduling algorithm that selects the process with the shortest burst time to execute next. It can be either preemptive or non-preemptive. In the non-preemptive version, once a process starts execution, it continues until completion without interruption.

Let's denote $n$ as the total number of processes, $AT_i$ as the arrival time of process $i$, $BT_i$ as the burst time of process $i$, and $CT_i$ as the completion time of process $i$. The completion time for each process can be calculated using the formula:

\[
CT_i = CT_{i-1} + BT_i
\]

where $CT_0 = 0$. The average waiting time ($\text{AWT}$) for all processes can be computed similarly to FCFS.

Now, let's provide the algorithm for the SJF technique and explain it in detail:
\FloatBarrier
\begin{algorithm}[H]
\caption{Shortest Job First (SJF) Scheduling}
\begin{algorithmic}[1]
\State Sort processes based on burst time $BT_i$
\State $CT_1 \gets AT_1 + BT_1$
\For{$i \gets 2$ to $n$}
    \State $CT_i \gets CT_{i-1} + BT_i$
\EndFor
\end{algorithmic}
\end{algorithm}
\FloatBarrier
The SJF algorithm sorts processes based on their burst time in ascending order and then follows the same steps as FCFS to calculate completion times.

\subsection{Priority Scheduling}

Priority Scheduling is a scheduling algorithm where each process is assigned a priority, and the process with the highest priority is selected for execution first. It can be either preemptive or non-preemptive.

Let's denote $n$ as the total number of processes, $AT_i$ as the arrival time of process $i$, $BT_i$ as the burst time of process $i$, and $CT_i$ as the completion time of process $i$. The completion time for each process can be calculated similarly to FCFS.

Now, let's provide the algorithm for the Priority Scheduling technique and explain it in detail:
\FloatBarrier
\begin{algorithm}[H]
\caption{Priority Scheduling}
\begin{algorithmic}[1]
\State Assign priority to each process
\State Sort processes based on priority
\State $CT_1 \gets AT_1 + BT_1$
\For{$i \gets 2$ to $n$}
    \State $CT_i \gets CT_{i-1} + BT_i$
\EndFor
\end{algorithmic}
\end{algorithm}
\FloatBarrier
The Priority Scheduling algorithm assigns priorities to processes and then sorts them based on their priority. It then calculates completion times similarly to FCFS.

\subsection{Round-Robin Scheduling}

Round-Robin Scheduling is a preemptive scheduling algorithm where each process is assigned a fixed time quantum, and processes are executed in a circular queue. If a process does not complete within its time quantum, it is preempted and placed at the end of the queue to wait for its turn again.

Let's denote $n$ as the total number of processes, $AT_i$ as the arrival time of process $i$, $BT_i$ as the burst time of process $i$, $CT_i$ as the completion time of process $i$, and $TQ$ as the time quantum.

The completion time for each process can be calculated using the formula:

\[
CT_i = \min(CT_{i-1} + TQ, AT_i + BT_i)
\]

where $CT_0 = 0$. The average waiting time ($\text{AWT}$) for all processes can be computed similarly to FCFS.

Now, let's provide the algorithm for the Round-Robin Scheduling technique and explain it in detail:
\FloatBarrier
\begin{algorithm}[H]
\caption{Round-Robin Scheduling}
\begin{algorithmic}[1]
\State Set $CT_1 \gets \min(AT_1 + BT_1, TQ)$
\For{$i \gets 2$ to $n$}
    \State $CT_i \gets \min(CT_{i-1} + TQ, AT_i + BT_i)$
\EndFor
\end{algorithmic}
\end{algorithm}
\FloatBarrier
The Round-Robin Scheduling algorithm assigns a fixed time quantum to each process and processes them in a circular queue, calculating completion times accordingly.


\section{Advanced Scheduling Techniques}

Scheduling techniques are essential in computer science and operating systems for efficiently managing resources and executing tasks. In this section, we delve into advanced scheduling techniques, including preemptive vs non-preemptive scheduling, multiprocessor scheduling, and real-time scheduling.

\subsection{Preemptive vs Non-Preemptive Scheduling}

Preemptive and non-preemptive scheduling are two fundamental approaches to task scheduling in operating systems. Preemptive scheduling allows the scheduler to interrupt a running task to allocate CPU time to another task with higher priority, while non-preemptive scheduling does not allow such interruptions.

\subsubsection{Preemptive Scheduling Algorithm}

Preemptive scheduling algorithms prioritize tasks based on their priority levels, executing higher-priority tasks before lower-priority ones. One such algorithm is the Shortest Remaining Time First (SRTF) algorithm, which selects the task with the shortest remaining processing time whenever a new task arrives or a running task completes.
\FloatBarrier
\begin{algorithm}
\caption{Shortest Remaining Time First (SRTF)}
\begin{algorithmic}[1]
\Procedure{SRTF}{tasks}
\State $current\_time \gets 0$
\While{$tasks$ not empty}
\State $next\_task \gets$ task with shortest remaining time
\State $execute(next\_task)$
\State $current\_time \gets current\_time + next\_task.duration$
\EndWhile
\EndProcedure
\end{algorithmic}
\end{algorithm}
\FloatBarrier
\subsubsection{Non-Preemptive Scheduling Algorithm}

Non-preemptive scheduling algorithms execute tasks until completion without interruption. One such algorithm is the Shortest Job First (SJF) algorithm, which selects the task with the shortest processing time among all waiting tasks.
\FloatBarrier
\begin{algorithm}
\caption{Shortest Job First (SJF)}
\begin{algorithmic}[1]
\Procedure{SJF}{tasks}
\State $current\_time \gets 0$
\While{$tasks$ not empty}
\State $next\_task \gets$ task with shortest processing time
\State $execute(next\_task)$
\State $current\_time \gets current\_time + next\_task.duration$
\EndWhile
\EndProcedure
\end{algorithmic}
\end{algorithm}

\subsubsection{Comparison of Preemptive vs Non-Preemptive Scheduling}

Preemptive and non-preemptive scheduling are two contrasting approaches to task scheduling, each with its advantages and disadvantages.

\textbf{Preemptive Scheduling:}
\begin{itemize}
    \item \textbf{Responsiveness:} Preemptive scheduling offers better responsiveness as higher-priority tasks can be executed immediately when they become runnable. This is particularly beneficial in time-critical systems where tasks need to respond quickly to external events.
    \item \textbf{Dynamic Priority Adjustment:} Preemptive scheduling allows for dynamic adjustment of task priorities based on changing system conditions. Tasks with higher priority can preempt lower-priority tasks, ensuring that critical tasks are executed promptly.
    \item \textbf{Overhead:} However, preemptive scheduling introduces overhead due to frequent context switches. When a higher-priority task preempts a lower-priority task, the processor must save the state of the preempted task and restore the state of the preempting task, leading to increased overhead.
\end{itemize}

\textbf{Non-Preemptive Scheduling:}
\begin{itemize}
    \item \textbf{Predictability:} Non-preemptive scheduling offers better predictability as tasks are executed until completion without interruption. This can simplify task management and analysis, especially in real-time systems where deterministic behavior is crucial.
    \item \textbf{Reduced Overhead:} Non-preemptive scheduling reduces overhead compared to preemptive scheduling since there are no context switches during task execution. This can lead to improved system efficiency, especially on systems with limited resources.
    \item \textbf{Response Time Variability:} However, non-preemptive scheduling may result in longer response times for high-priority tasks if they are blocked by lower-priority tasks. Tasks with higher priority must wait for lower-priority tasks to complete before they can be executed, leading to increased variability in response times.
\end{itemize}

\textbf{Choosing Between Preemptive and Non-Preemptive Scheduling:}
The choice between preemptive and non-preemptive scheduling depends on the specific requirements of the system and the characteristics of the tasks being scheduled. In time-critical systems where responsiveness is paramount, preemptive scheduling may be preferred despite the overhead. In contrast, for systems where predictability and determinism are more important, non-preemptive scheduling may be a better choice to ensure consistent behavior.

In practice, hybrid approaches that combine preemptive and non-preemptive scheduling techniques may be used to balance responsiveness and predictability, depending on the system's requirements and constraints.


\subsection{Multiprocessor Scheduling}

Multiprocessor scheduling involves efficiently allocating tasks to multiple processors for parallel execution, aiming to maximize throughput and minimize resource contention.

\subsubsection{Multiprocessor Scheduling Algorithm}

One approach to multiprocessor scheduling is the Work-Conserving algorithm, which ensures that each processor remains busy as much as possible by assigning tasks dynamically based on their computational requirements.
\FloatBarrier
\begin{algorithm}
\caption{Work-Conserving Multiprocessor Scheduling}
\begin{algorithmic}[1]
\Procedure{Work-Conserving}{tasks, processors}
\While{$tasks$ not empty}
\State $task \gets$ select task with highest priority
\State $processor \gets$ select idle processor
\State $assign(task, processor)$
\EndWhile
\EndProcedure
\end{algorithmic}
\end{algorithm}

\subsection{Real-Time Scheduling}

Real-time scheduling is crucial for systems where tasks must meet strict deadlines to ensure proper operation, such as in embedded systems and multimedia applications.

\subsubsection{Real-Time Scheduling Algorithm}

The Rate-Monotonic Scheduling (RMS) algorithm assigns priorities to tasks based on their periods, with shorter periods assigned higher priorities. This ensures that tasks with tighter deadlines are executed before those with looser deadlines.
\FloatBarrier
\begin{algorithm}
\caption{Rate-Monotonic Scheduling (RMS)}
\begin{algorithmic}[1]
\Procedure{RMS}{tasks}
\State $current\_time \gets 0$
\While{$tasks$ not empty}
\State $next\_task \gets$ task with highest priority
\State $execute(next\_task)$
\State $current\_time \gets current\_time + next\_task.period$
\EndWhile
\EndProcedure
\end{algorithmic}
\end{algorithm}

\section{Scheduling in Operating Systems}
Operating systems employ various scheduling techniques to manage and optimize the allocation of resources such as CPU time, memory, and I/O devices. Process scheduling involves determining the order in which processes are executed by the CPU. Thread scheduling focuses on managing the execution of individual threads within a process. Additionally, I/O scheduling aims to optimize the order in which I/O requests are serviced by devices such as disks and network interfaces.

\subsection{Process Scheduling}
Process scheduling is a fundamental aspect of operating system design, as it directly impacts system performance and responsiveness. The goal of process scheduling is to maximize system throughput, minimize response time, and ensure fair allocation of resources among competing processes.

One of the most commonly used process scheduling algorithms is the \textbf{Round Robin (RR)} scheduling algorithm. RR is a preemptive scheduling algorithm that assigns a fixed time quantum to each process in the ready queue. When a process's time quantum expires, it is preempted and moved to the end of the queue, allowing the next process to execute.

Let $T_i$ represent the total execution time of process $i$, and $t_i$ represent the time quantum assigned to process $i$ in RR scheduling. The average waiting time ($W_{\text{avg}}$) for all processes in the queue can be calculated using the following formula:

\[
W_{\text{avg}} = \frac{\sum_{i=1}^{n} (W_i)}{n}
\]

Where $W_i$ represents the waiting time for process $i$, which can be calculated as:

\[
W_i = (n - 1) \times t_i
\]

The RR scheduling algorithm ensures fairness by providing each process with an equal share of CPU time, regardless of its priority or execution characteristics.
\FloatBarrier
\begin{algorithm}
\caption{Round Robin (RR) Scheduling Algorithm}\label{rr}
\begin{algorithmic}[1]
\Procedure{RoundRobin}{$processes, quantum$}
\State $queue \gets empty$
\State $time \gets 0$
\While{$\text{not } processes.empty()$}
\State $process \gets processes.pop()$
\State $execute(process, quantum)$
\State $time \gets time + quantum$
\If{$process \text{ not finished}$}
\State $queue.push(process)$
\EndIf
\EndWhile
\EndProcedure
\end{algorithmic}
\end{algorithm}
\FloatBarrier
\textbf{Python Code Implementation:}
\FloatBarrier
\begin{lstlisting}
def round_robin(processes, quantum):
    queue = []
    time = 0
    while processes:
        process = processes.pop(0)
        execute(process, quantum)
        time += quantum
        if not process.finished:
            queue.append(process)
    return
\end{lstlisting}
\FloatBarrier

\subsection{Thread Scheduling}
Thread scheduling focuses on managing the execution of individual threads within a process. One popular thread scheduling algorithm is the \textbf{Multilevel Feedback Queue (MLFQ)} algorithm. MLFQ maintains multiple queues with different priority levels, where threads are scheduled based on their priority and recent execution history.
\FloatBarrier
\begin{algorithm}
\caption{Multilevel Feedback Queue (MLFQ) Scheduling Algorithm}\label{mlfq}
\begin{algorithmic}[1]
\Procedure{MLFQ}{$threads$}
\State $queues \gets initialize\_queues()$
\While{$\text{not all queues are empty}$}
\For{$i \gets 0$ to $n$}
\If{$queue[i] \text{ not empty}$}
\State $thread \gets queue[i].pop()$
\State $execute(thread)$
\If{$thread \text{ not finished}$}
\State $queues.adjust\_priority(thread)$
\EndIf
\EndIf
\EndFor
\EndWhile
\EndProcedure
\end{algorithmic}
\end{algorithm}
\FloatBarrier
\textbf{Python Code Implementation:}
\FloatBarrier
\begin{lstlisting}
def MLFQ(threads):
    queues = initialize_queues()
    while not all_empty(queues):
        for queue in queues:
            if queue:
                thread = queue.pop(0)
                execute(thread)
                if not thread.finished:
                    queues.adjust_priority(thread)
    return
\end{lstlisting}

\subsection{I/O Scheduling}
I/O scheduling is concerned with optimizing the order in which I/O requests are serviced by devices such as disks and network interfaces. One commonly used I/O scheduling algorithm is the \textbf{Shortest Seek Time First (SSTF)} algorithm. SSTF selects the I/O request that is closest to the current disk head position, minimizing the seek time required to access data.
\FloatBarrier
\begin{algorithm}
\caption{Shortest Seek Time First (SSTF) Scheduling Algorithm}\label{sstf}
\begin{algorithmic}[1]
\Procedure{SSTF}{$requests, head\_position$}
\State $sorted\_requests \gets sort\_by\_distance(requests, head\_position)$
\For{$request$ in $sorted\_requests$}
\State $execute(request)$
\EndFor
\EndProcedure
\end{algorithmic}
\end{algorithm}
\FloatBarrier
\textbf{Python Code Implementation:}
\FloatBarrier
\begin{lstlisting}
def SSTF(requests, head_position):
    sorted_requests = sort_by_distance(requests, head_position)
    for request in sorted_requests:
        execute(request)
    return
\end{lstlisting}


\section{Scheduling in Production and Manufacturing}
Scheduling plays a crucial role in production and manufacturing processes by determining the sequence of tasks or jobs to be executed on machines or resources over time. Effective scheduling leads to improved efficiency, reduced production time, and optimized resource utilization. In this section, we will explore three important scheduling techniques: Job Shop Scheduling, Flow Shop Scheduling, and Project Scheduling.

\subsection{Job Shop Scheduling}
Job Shop Scheduling involves scheduling a set of jobs that need to be processed on a set of machines, where each job consists of multiple operations and each operation must be processed on a specific machine in a predefined order. This problem arises in various manufacturing settings, such as automobile assembly lines and semiconductor fabrication plants.

One of the most common algorithms used for Job Shop Scheduling is the \textbf{Johnson's Rule}. This algorithm aims to minimize the total completion time of all jobs by sequencing the operations optimally. Johnson's Rule works by identifying the critical path, which is the sequence of operations that determines the overall completion time, and then scheduling the operations accordingly.

\textbf{Algorithm: Johnson's Rule}
\begin{enumerate}
    \item For each job, compute the total processing time as the sum of processing times for all operations.
    \item Initialize two queues: one for machines with the smallest processing time first (Queue 1) and the other for machines with the largest processing time first (Queue 2).
    \item While both queues are not empty:
    \begin{itemize}
        \item Choose the operation with the smallest processing time from Queue 1 and schedule it next.
        \item Choose the operation with the smallest processing time from Queue 2 and schedule it next.
    \end{itemize}
\end{enumerate}

Let's illustrate Johnson's Rule with an example:
\FloatBarrier
\FloatBarrier
\begin{algorithm}
\caption{Johnson's Rule}
\begin{algorithmic}[1]
\Procedure{Johnson}{Jobs, Machines}
\State Compute total processing time for each job
\State Initialize Queues 1 and 2
\While{Both Queues are not empty}
    \State Schedule operation with smallest processing time from Queue 1
    \State Schedule operation with smallest processing time from Queue 2
\EndWhile
\EndProcedure
\end{algorithmic}
\end{algorithm}
\FloatBarrier

\subsection{Flow Shop Scheduling}
Flow Shop Scheduling involves scheduling a set of jobs that need to be processed on a series of machines in the same order, where each job consists of multiple operations and each operation must be processed on a specific machine in a predefined order. This problem arises in manufacturing settings where the sequence of operations is fixed, such as in assembly lines.

One of the most common algorithms used for Flow Shop Scheduling is the \textbf{Johnson's Algorithm}. This algorithm extends Johnson's Rule to the flow shop setting by scheduling jobs on two machines in a way that minimizes the total completion time.

\textbf{Algorithm: Johnson's Algorithm}
\begin{enumerate}
    \item For each job, compute the total processing time as the sum of processing times for all operations.
    \item Initialize two queues: one for machines with the smallest processing time first (Queue 1) and the other for machines with the largest processing time first (Queue 2).
    \item While both queues are not empty:
    \begin{itemize}
        \item Choose the operation with the smallest processing time from Queue 1 and schedule it next.
        \item Choose the operation with the smallest processing time from Queue 2 and schedule it next.
    \end{itemize}
\end{enumerate}

\subsection{Project Scheduling}
Project Scheduling involves scheduling a set of activities or tasks that need to be completed within a specified time frame, taking into account dependencies between tasks and resource constraints. This problem arises in project management, construction projects, and software development.

One of the most common algorithms used for Project Scheduling is the \textbf{Critical Path Method (CPM)}. CPM aims to determine the longest path of dependent activities, known as the critical path, and the total duration of the project. By identifying the critical path, project managers can allocate resources efficiently and ensure timely project completion.

\textbf{Algorithm: Critical Path Method (CPM)}
\begin{enumerate}
    \item Construct a network diagram representing all tasks and their dependencies.
    \item Compute the earliest start time (ES) and earliest finish time (EF) for each task.
    \item Compute the latest start time (LS) and latest finish time (LF) for each task.
    \item Calculate the slack time for each task (slack = LF - EF).
    \item Identify the critical path, which consists of tasks with zero slack time.
    \item Determine the total duration of the project as the sum of durations along the critical path.
\end{enumerate}

\section{Scheduling in Distributed Systems}

Scheduling in distributed systems involves the allocation of tasks to resources in an efficient and balanced manner to optimize performance and resource utilization. In this section, we will explore various aspects of scheduling in distributed systems, including load balancing, task allocation, and cloud computing scheduling.

\subsection{Load Balancing}

Load balancing is a critical aspect of distributed systems where the workload needs to be evenly distributed across multiple resources to ensure optimal resource utilization, prevent overload, and improve system performance. In this subsection, we will delve deeper into the concept of load balancing and explore additional load balancing algorithms.

Load balancing algorithms aim to distribute tasks or requests among available resources in a way that minimizes response time, maximizes throughput, and ensures fairness. These algorithms consider various factors such as resource capacity, current load, task characteristics, and system topology to make intelligent allocation decisions.

\subsubsection{Need for Load Balancing}

In distributed systems, the workload is often dynamic and unpredictable, with fluctuations in demand and resource availability. Without effective load balancing, some resources may become overwhelmed with tasks, leading to performance degradation and potential system failures, while others may remain underutilized. Load balancing helps address these challenges by dynamically redistributing tasks to maintain a balanced workload across the system

\subsubsection{Least Loaded Algorithm}

The Least Loaded Algorithm assigns incoming tasks to the resource with the least load, where load can be measured in terms of CPU utilization, memory usage, or other relevant metrics. Mathematically, the algorithm can be described as follows:

Let $R$ be the set of resources, and $L(r)$ be the load on resource $r \in R$. When a new task $T$ arrives, the algorithm selects the resource $r^*$ such that:

$$r^* = \text{argmin}_{r \in R} L(r)$$

The task $T$ is then assigned to the resource $r^*$.
\FloatBarrier
\begin{algorithm}[H]
\caption{Least Loaded Algorithm}\label{alg:least_loaded}
\begin{algorithmic}[1]
\Procedure{LeastLoaded}{$R, T$}
\State $r^* \gets \text{argmin}_{r \in R} L(r)$
\State Assign task $T$ to resource $r^*$
\EndProcedure
\end{algorithmic}
\end{algorithm}
\FloatBarrier   
\textbf{Python Code Implementation:}
\FloatBarrier
\begin{lstlisting}
def least_loaded(resources, task):
    least_loaded_resource = min(resources, key=lambda r: r.load)
    least_loaded_resource.assign_task(task)
\end{lstlisting}
\FloatBarrier
The Least Loaded Algorithm selects the resource with the minimum load and assigns the incoming task to that resource. This helps in achieving load balancing across distributed systems efficiently.

\subsection{Task Allocation}

Task allocation is a critical aspect of distributed systems where tasks need to be assigned to resources in a way that optimizes system performance, resource utilization, and overall efficiency. In this subsection, we will delve deeper into the concept of task allocation and explore additional task allocation algorithms.

Task allocation algorithms aim to assign tasks to resources based on various factors such as task characteristics, resource capabilities, and system constraints. These algorithms help optimize resource utilization, minimize task completion time, and ensure fairness in task distribution.

\subsubsection{Need for Task Allocation}

In distributed systems, tasks may have different requirements and constraints, and resources may vary in terms of capacity, performance, and availability. Effective task allocation is essential to ensure that tasks are executed on appropriate resources to meet performance objectives, resource constraints, and service-level agreements.

\subsubsection{Task Allocation Algorithm}

One common approach to task allocation is the Max-Min Algorithm, which aims to maximize the minimum resource allocation for each task.
\FloatBarrier
\begin{algorithm}[H]
\caption{Max-Min Algorithm}\label{alg:max_min}
\begin{algorithmic}[1]
\Procedure{MaxMin}{$T, R$}
\State Sort tasks in descending order of resource requirements
\State Sort resources in ascending order of available capacity
\For{each task $t$}
    \State Allocate $t$ to the resource with the most available capacity
\EndFor
\EndProcedure
\end{algorithmic}
\end{algorithm}
\FloatBarrier   
\textbf{Python Code Implementation:}
\FloatBarrier
\begin{lstlisting}
def max_min(tasks, resources):
    tasks.sort(key=lambda t: t.requirements, reverse=True)
    resources.sort(key=lambda r: r.available_capacity)
    
    for task in tasks:
        resource = max(resources, key=lambda r: r.available_capacity)
        resource.assign_task(task)
\end{lstlisting}
\FloatBarrier
The Max-Min Algorithm optimizes task allocation by maximizing the minimum resource allocation for each task, thereby improving system performance and resource utilization.

\subsection{Cloud Computing Scheduling}

Cloud computing scheduling involves efficiently allocating tasks to virtual machines (VMs) or containers in cloud environments to meet service-level objectives and minimize costs. In this subsection, we will delve deeper into the concept of cloud computing scheduling and explore additional scheduling algorithms.

\subsubsection{Need for Cloud Computing Scheduling}

Cloud computing environments typically host a large number of VMs or containers, each running various tasks or services. Efficiently allocating tasks to these resources is essential to maximize resource utilization, ensure performance scalability, and optimize cost-effectiveness. Cloud computing scheduling helps address these challenges by dynamically assigning tasks to resources based on workload characteristics, resource availability, and system requirements.


\subsubsection{Cloud Computing Scheduling Algorithm}

One widely used scheduling algorithm in cloud computing is the First Fit Decreasing (FFD) algorithm, which sorts tasks in descending order of resource requirements and allocates them to the first VM or container that has sufficient capacity.
\FloatBarrier
\begin{algorithm}[H]
\caption{First Fit Decreasing (FFD) Algorithm}\label{alg:ffd}
\begin{algorithmic}[1]
\Procedure{FFD}{$T, VMs$}
\State Sort tasks in descending order of resource requirements
\For{each task $t$}
    \For{each VM $v$}
        \If{$t$ fits in $v$}
            \State Allocate $t$ to $v$
            \State Break
        \EndIf
    \EndFor
\EndFor
\EndProcedure
\end{algorithmic}
\end{algorithm}
\FloatBarrier
\textbf{Python Code Implementation:}
\FloatBarrier
\begin{lstlisting}
def ffd(tasks, vms):
    tasks.sort(key=lambda t: t.requirements, reverse=True)
    
    for task in tasks:
        for vm in vms:
            if task.fits(vm):
                vm.assign_task(task)
                break
\end{lstlisting}
\FloatBarrier
The FFD algorithm efficiently allocates tasks to VMs in cloud environments, helping optimize resource utilization and meet service-level objectives effectively.


\section{Scheduling in Networks}

Scheduling in networks plays a critical role in optimizing the allocation of resources and ensuring efficient communication. This section explores various aspects of scheduling techniques in network systems, including packet scheduling in network routers, bandwidth allocation, and quality of service (QoS) management.

\subsection{Packet Scheduling in Network Routers}

Packet scheduling in network routers is a fundamental aspect of traffic management, where incoming packets are scheduled for transmission based on certain criteria. This ensures fair distribution of resources and minimizes congestion in the network.

Packet scheduling algorithms aim to allocate resources efficiently while meeting certain performance criteria such as delay, throughput, and fairness. One widely used packet scheduling algorithm is Weighted Fair Queuing (WFQ).

\subsubsection{Mathematical Discussion}

Let $N$ be the number of flows in the network, and let $w_i$ denote the weight assigned to flow $i$, where $1 \leq i \leq N$. The weight represents the relative priority or importance of each flow. Let $R_i$ be the rate allocated to flow $i$, representing the amount of bandwidth assigned to the flow.

The WFQ algorithm assigns a virtual finish time $F_i$ to each packet in flow $i$, calculated as:

\begin{equation*}
    F_i = \text{Arrival time} + \frac{\text{Packet size}}{R_i}
\end{equation*}

The packets are then enqueued in a priority queue based on their virtual finish times, and transmitted in the order of increasing virtual finish times.

\subsubsection{Algorithmic Example: Weighted Fair Queuing (WFQ)}
\FloatBarrier
\begin{algorithm}[H]
\caption{Weighted Fair Queuing (WFQ)}
\begin{algorithmic}[1]
\State Initialize virtual time for each flow
\State Initialize empty priority queue
\While{packets arrive}
    \State Assign virtual finish time for each packet based on its arrival time and service rate
    \State Enqueue packets into priority queue based on virtual finish time
    \State Dequeue and transmit packets in order of increasing virtual finish time
\EndWhile
\end{algorithmic}
\end{algorithm}
\FloatBarrier

\subsection{Bandwidth Allocation}

Bandwidth allocation involves dividing the available bandwidth among competing users or applications in a network. This ensures that each user receives a fair share of the available resources while maximizing overall network efficiency.

\subsubsection{Mathematical Discussion}

In bandwidth allocation, the goal is to distribute the available bandwidth among $N$ users or applications in a fair and efficient manner. Let $X_i$ denote the bandwidth allocated to user $i$, where $1 \leq i \leq N$. The total available bandwidth is denoted by $B$.

One common approach to bandwidth allocation is Max-Min Fairness. In this approach, the bandwidth allocation vector $\mathbf{X}$ is iteratively adjusted until convergence. At each iteration, the deficit for each user is calculated based on the current allocation. The deficit represents the difference between the user's allocated bandwidth and its fair share. Bandwidth is then allocated to each user proportional to its deficit until convergence is achieved.

\subsubsection{Algorithmic Example: Max-Min Fairness}
\FloatBarrier
\begin{algorithm}[H]
\caption{Max-Min Fairness}
\begin{algorithmic}[1]
\State Initialize bandwidth allocation vector $\mathbf{X}$ with equal shares
\While{not converged}
    \State Calculate the deficit for each user based on current allocation
    \State Allocate bandwidth to each user proportional to its deficit until convergence
\EndWhile
\end{algorithmic}
\end{algorithm}
\FloatBarrier

\subsection{Quality of Service (QoS) Management}

Quality of Service (QoS) management involves prioritizing certain types of traffic or users to meet specific performance requirements, such as latency, throughput, and reliability. This ensures that critical applications receive the necessary resources to function optimally.

\subsubsection{Mathematical Discussion}

In QoS management, different traffic classes are assigned different priority levels based on their QoS requirements. Let $T_i$ denote the traffic class of flow $i$, where $1 \leq i \leq N$. Each traffic class is associated with a set of performance parameters, such as latency bounds, packet loss rates, and throughput requirements.

One approach to QoS management is the Token Bucket Algorithm. In this algorithm, a token bucket is used to control the rate at which packets are transmitted. Each packet is allowed to be transmitted only if there are sufficient tokens in the bucket. If tokens are depleted, packets are either dropped or queued for later transmission.

\subsubsection{Algorithmic Example: Token Bucket Algorithm}
\FloatBarrier
\begin{algorithm}[H]
\caption{Token Bucket Algorithm}
\begin{algorithmic}[1]
\State Initialize token bucket with tokens and token replenishment rate
\While{packets arrive}
    \If{tokens available}
        \State Transmit packet
        \State Consume tokens
    \Else
        \State Drop packet
    \EndIf
\EndWhile
\end{algorithmic}
\end{algorithm}
\FloatBarrier

\section{Mathematical and Heuristic Methods in Scheduling}

Scheduling problems involve allocating limited resources over time to perform a set of tasks, subject to various constraints and objectives. Mathematical and heuristic methods are commonly used to solve scheduling problems efficiently and effectively. In this section, we explore different approaches, including Linear Programming, Heuristic Algorithms, and Evolutionary Algorithms.

\subsection{Linear Programming Approaches}

Linear Programming (LP) provides a mathematical framework for optimizing linear objective functions subject to linear constraints. In the context of scheduling, LP approaches formulate the scheduling problem as a linear optimization problem and solve it using LP techniques.

\begin{itemize}
    \item \textbf{Job Scheduling with Minimization of Total Completion Time:} In this approach, the objective is to minimize the total completion time of all jobs subject to various constraints such as job dependencies and resource availability. Mathematically, this can be formulated as follows:
    \FloatBarrier
    \begin{align*}
        & \text{Minimize} \quad \sum_{i=1}^{n} C_i \\
        & \text{Subject to} \\
        & C_i = \text{start time of job } i \\
        & C_i + p_i \leq C_j \quad \forall (i, j) \in \text{precedence constraints} \\
        & C_i + p_i \leq d_i \quad \forall i \in \text{jobs} \\
        & C_i \geq 0 \quad \forall i \in \text{jobs}
    \end{align*}
    \FloatBarrier
    where \(C_i\) represents the completion time of job \(i\), \(p_i\) is the processing time of job \(i\), and \(d_i\) is the due date of job \(i\). The first constraint ensures that each job starts after its predecessor completes, the second constraint enforces the due date constraint, and the last constraint ensures non-negative start times.
    
    \item \textbf{Job Scheduling with Minimization of Total Weighted Tardiness:} This approach aims to minimize the total weighted tardiness of jobs, where tardiness is the amount of time a job finishes after its due date, weighted by a given weight for each job. Mathematically, this can be formulated as follows:
    \FloatBarrier
    \begin{align*}
        & \text{Minimize} \quad \sum_{i=1}^{n} w_i \cdot \max(0, C_i - d_i) \\
        & \text{Subject to} \\
        & C_i = \text{start time of job } i \\
        & C_i + p_i \leq C_j \quad \forall (i, j) \in \text{precedence constraints} \\
        & C_i + p_i \leq d_i \quad \forall i \in \text{jobs} \\
        & C_i \geq 0 \quad \forall i \in \text{jobs}
    \end{align*}
    \FloatBarrier
    where \(w_i\) is the weight of job \(i\), and all other constraints are similar to the previous formulation.
\end{itemize}

Let's illustrate the Job Scheduling with Minimization of Total Completion Time using the following algorithmic example:
\FloatBarrier
\begin{algorithm}
\caption{LP-based Job Scheduling with Minimization of Total Completion Time}
\begin{algorithmic}[1]
\State Formulate the LP problem as described above.
\State Solve the LP problem using an LP solver (e.g., simplex method).
\State Extract the optimal solution to obtain the start times of each job.
\end{algorithmic}
\end{algorithm}
\FloatBarrier

\subsection{Heuristic Algorithms}

Heuristic algorithms provide approximate solutions to scheduling problems in a reasonable amount of time. These algorithms do not guarantee optimality but often produce high-quality solutions efficiently.

\begin{itemize}
    \item \textbf{Earliest Due Date (EDD) Algorithm:} The EDD algorithm schedules jobs based on their due dates, with the earliest due date job scheduled first. Mathematically, the EDD rule can be expressed as:
    \[ \text{Priority}(i) = d_i \]
    where \(d_i\) is the due date of job \(i\).
    
    \item \textbf{Shortest Processing Time (SPT) Algorithm:} The SPT algorithm schedules jobs based on their processing times, with the shortest processing time job scheduled first. Mathematically, the SPT rule can be expressed as:
    \[ \text{Priority}(i) = p_i \]
    where \(p_i\) is the processing time of job \(i\).
\end{itemize}

Let's demonstrate the Earliest Due Date (EDD) Algorithm using the following algorithmic example:
\FloatBarrier
\begin{algorithm}
\caption{Earliest Due Date (EDD) Algorithm}
\begin{algorithmic}[1]
\For{each job \(i\) in non-decreasing order of due dates}
\State Schedule job \(i\) at the earliest available time slot.
\EndFor
\end{algorithmic}
\end{algorithm}
\FloatBarrier

\subsection{Evolutionary Algorithms and Machine Learning}

Evolutionary algorithms (EAs) are population-based optimization techniques inspired by the process of natural selection. These algorithms iteratively evolve a population of candidate solutions to search for optimal or near-optimal solutions to a given problem. In the context of scheduling, EAs are used to explore the solution space efficiently and find high-quality schedules.

\begin{itemize}
    \item \textbf{Genetic Algorithm (GA):} GA is one of the most well-known EAs. It operates on a population of potential solutions, where each solution is represented as a chromosome. The chromosome encodes a candidate solution to the scheduling problem, typically representing a permutation of job orders or a binary encoding of job allocations. GA iteratively applies genetic operators such as selection, crossover, and mutation to evolve better solutions over generations.

    \item \textbf{Particle Swarm Optimization (PSO):} PSO is another popular EA based on the social behavior of particles in a swarm. In PSO, each potential solution is represented as a particle moving in the search space. The position of each particle corresponds to a candidate solution, and the velocity of each particle determines how it explores the search space. PSO iteratively updates the position and velocity of particles based on their own best-known position and the global best-known position in the search space.

\end{itemize}

Let's provide more detailed algorithmic examples for both GA and PSO in the context of scheduling:

\subsubsection{Genetic Algorithm (GA)}

In the context of scheduling, a chromosome in GA represents a permutation of job orders. The objective is to find the best permutation that minimizes a scheduling criterion such as total completion time or total weighted tardiness.
\FloatBarrier
\begin{algorithm}
\caption{Genetic Algorithm for Job Scheduling}
\begin{algorithmic}[1]
\State Initialize population of chromosomes randomly.
\For{each generation}
    \State Evaluate fitness of each chromosome.
    \State Select parents for reproduction based on fitness.
    \State Apply crossover and mutation operators to create offspring.
    \State Replace old population with new population of offspring.
\EndFor
\Return Best chromosome found.
\end{algorithmic}
\end{algorithm}

\subsubsection{Particle Swarm Optimization (PSO)}

In PSO, each particle represents a potential solution, and its position corresponds to a candidate schedule. The objective is to find the best position (schedule) that minimizes a scheduling criterion.
\FloatBarrier
\begin{algorithm}
\caption{Particle Swarm Optimization for Job Scheduling}
\begin{algorithmic}[1]
\State Initialize particles with random positions and velocities.
\For{each iteration}
    \State Update personal best-known position for each particle.
    \State Update global best-known position for the swarm.
    \State Update particle positions and velocities.
\EndFor
\Return Best position found.
\end{algorithmic}
\end{algorithm}
\FloatBarrier
Both GA and PSO can be customized and extended to handle various scheduling criteria and constraints. They offer the advantage of exploring the solution space effectively and finding high-quality schedules, although they may not guarantee optimality.

Evolutionary algorithms are relevant to machine learning in scheduling as they can adaptively learn from past schedules and improve scheduling policies over time, leading to more efficient and effective scheduling solutions.


\section{Challenges and Future Directions}

Scheduling plays a crucial role in various domains, including manufacturing, transportation, project management, and computer systems. However, as systems become more complex and dynamic, new challenges and opportunities arise in the field of scheduling. In this section, we discuss some of the key challenges and future directions in scheduling research.

\subsection{Scalability Issues}

Scalability is a critical concern in scheduling algorithms, especially as systems grow larger and more complex. One of the main challenges is the computational complexity of scheduling problems, which often grows exponentially with the size of the problem instance. Let's formalize this concept using mathematical notation.

Consider a scheduling problem with \(n\) tasks and \(m\) resources. Let \(T\) denote the set of tasks \(T = \{1, 2, ..., n\}\), and let \(R\) denote the set of resources \(R = \{1, 2, ..., m\}\). Each task \(i \in T\) requires a certain amount of each resource \(j \in R\), denoted by \(r_{ij}\).

To represent the schedule, we introduce binary decision variables \(x_{ijt}\), where \(x_{ijt} = 1\) if task \(i\) is assigned to resource \(j\) at time \(t\), and \(x_{ijt} = 0\) otherwise. Additionally, let \(C_{ij}\) represent the capacity of resource \(j\) at time \(t\).

The scheduling problem can be formulated as an optimization problem, where the objective is to minimize or maximize some objective function subject to constraints. For example, the objective function could be to minimize the makespan, i.e., the total time required to complete all tasks, subject to constraints such as resource capacities and task dependencies.

\textbf{Mathematically Detailed Discussion:}

The computational complexity of scheduling problems depends on various factors, including the size of the problem instance, the complexity of task dependencies, and the presence of constraints.

One approach to modeling the computational complexity of scheduling problems is through complexity classes. For example, the class NP-hard encompasses problems for which no polynomial-time algorithm exists to solve them, but solutions can be verified in polynomial time.

Many scheduling problems are NP-hard, meaning that finding an optimal solution is computationally intractable in the general case. For example, the classical job shop scheduling problem, where each task must be processed on a specific machine in a specific order, is NP-hard.

To address scalability issues in scheduling algorithms, researchers often develop heuristic or approximation algorithms that provide near-optimal solutions in a reasonable amount of time. These algorithms sacrifice optimality for efficiency and are often based on greedy algorithms, metaheuristic techniques, or integer programming formulations.

In summary, scalability issues in scheduling algorithms stem from the exponential growth in computational complexity with the size of the problem instance. Addressing these issues requires the development of efficient algorithms and optimization techniques tailored to specific problem domains and constraints.


\subsection{Scheduling under Uncertainty}

Scheduling under uncertainty is an important area of research that deals with the challenges of making scheduling decisions in the presence of uncertain information. This uncertainty can arise from various sources, including variability in task durations, resource availability, and external factors such as weather conditions or market fluctuations.

Mathematically, scheduling under uncertainty can be modeled using stochastic optimization techniques, probabilistic models, and decision theory. One common approach is to formulate scheduling problems as stochastic optimization problems and use techniques such as stochastic programming, Markov decision processes, or robust optimization to find robust schedules that perform well under different scenarios.

To illustrate, let's consider a simple scheduling problem where tasks have uncertain durations. We can model this problem as a stochastic scheduling problem and use dynamic programming to find an optimal schedule that minimizes the expected completion time or maximizes resource utilization while considering the uncertainty in task durations.

\subsection{Emerging Applications and Technologies}

Scheduling techniques are finding new applications and opportunities in emerging domains and technologies. Some of the emerging applications and technologies driving research in scheduling include:

\begin{itemize}
    \item \textbf{Smart Manufacturing}: Scheduling algorithms are being applied to optimize production processes in smart manufacturing systems, where machines, sensors, and devices are interconnected to improve efficiency, reduce downtime, and minimize energy consumption.
    
    \item \textbf{Internet of Things (IoT)}: Scheduling algorithms are used to coordinate and manage the resources and tasks in IoT networks, where a large number of devices with limited computational and communication capabilities need to be scheduled efficiently.
    
    \item \textbf{Blockchain Technology}: Scheduling algorithms play a role in scheduling transactions and smart contracts in blockchain networks, where decentralized and distributed consensus mechanisms require efficient resource allocation and task scheduling.
    
    \item \textbf{Edge Computing}: Scheduling algorithms are applied to optimize task scheduling and resource allocation in edge computing environments, where computation and storage resources are distributed closer to the end-users to reduce latency and improve performance.
    
    \item \textbf{Autonomous Systems}: Scheduling algorithms are used in autonomous systems such as autonomous vehicles, drones, and robots to plan and schedule tasks, optimize routes, and allocate resources autonomously in dynamic and uncertain environments.
\end{itemize}

Each of these emerging applications presents unique challenges and opportunities for scheduling research, driving the development of new algorithmic techniques and optimization methods tailored to specific domains and technologies.

\section{Case Studies}
Scheduling techniques play a crucial role in various domains, including large-scale data centers, smart grids, and autonomous systems. In this section, we delve into the intricacies of scheduling in these contexts, exploring the challenges, algorithms, and applications associated with each.

\subsection{Scheduling in Large-Scale Data Centers}
Large-scale data centers are key components of modern computing infrastructure, hosting a myriad of applications and services. Efficient scheduling in data centers is essential for optimizing resource utilization, minimizing latency, and maximizing throughput. In this subsection, we examine the mathematical foundations and algorithmic techniques used for scheduling tasks in large-scale data centers.

\subsubsection{Mathematically Detailed Discussion}
In a large-scale data center, tasks arrive dynamically and need to be assigned to available resources such as servers or virtual machines (VMs). The scheduling problem can be formulated as an optimization task where the objective is to minimize the total completion time or maximize resource utilization while satisfying constraints such as deadlines or service-level agreements.

One commonly used algorithm for scheduling tasks in large-scale data centers is the \textbf{Minimum Completion Time (MCT)} algorithm. The MCT algorithm assigns each task to the resource that minimizes its completion time, considering factors such as processing power, network bandwidth, and data locality. The algorithm can be described as follows:
\FloatBarrier
\begin{algorithm}[H]
\caption{Minimum Completion Time (MCT) Algorithm}
\begin{algorithmic}[1]
\State Sort tasks in non-decreasing order of processing time
\For{each task $T$}
\State Assign $T$ to the resource that minimizes its completion time
\EndFor
\end{algorithmic}
\end{algorithm}

\subsubsection{Algorithmic Example}
Here's a Python implementation of the MCT algorithm for scheduling tasks in a large-scale data center:
\FloatBarrier
\begin{lstlisting}
def mct_algorithm(tasks, resources):
    tasks.sort(key=lambda x: x.processing_time)
    for task in tasks:
        min_completion_time = float('inf')
        selected_resource = None
        for resource in resources:
            completion_time = calculate_completion_time(task, resource)
            if completion_time < min_completion_time:
                min_completion_time = completion_time
                selected_resource = resource
        assign_task_to_resource(task, selected_resource)
\end{lstlisting}
\FloatBarrier

\subsection{Dynamic Scheduling in Smart Grids}
Smart grids leverage advanced communication and control technologies to efficiently manage electricity generation, distribution, and consumption. Dynamic scheduling techniques are vital for optimizing energy usage, minimizing costs, and maintaining grid stability. In this subsection, we explore the mathematical models and algorithmic approaches used for dynamic scheduling in smart grids.

\subsubsection{Mathematically Detailed Discussion}
In a smart grid, dynamic scheduling involves real-time allocation of energy resources such as power generators, storage systems, and demand-side resources to meet electricity demand while adhering to operational constraints and market dynamics. The scheduling problem can be formulated as an optimization task aiming to minimize energy costs, reduce carbon emissions, or enhance grid reliability.

One prevalent technique for dynamic scheduling in smart grids is \textbf{Model Predictive Control (MPC)}. MPC formulates the scheduling problem as a dynamic optimization task and solves it iteratively based on predictions of future system behavior. The algorithm adjusts resource allocations in response to changing conditions, such as fluctuating demand or renewable energy availability.

\subsubsection{Algorithmic Example}
Here's a simplified Python implementation of the Model Predictive Control (MPC) algorithm for dynamic scheduling in a smart grid:
\FloatBarrier
\begin{lstlisting}
def mpc_algorithm():
    initialize_system_state()
    for t in range(total_time_steps):
        predict_future_state()
        optimize_resource_allocation()
        update_system_state()
\end{lstlisting}
\FloatBarrier

\subsection{Adaptive Scheduling in Autonomous Systems}
Autonomous systems such as unmanned aerial vehicles (UAVs) and self-driving cars rely on adaptive scheduling techniques to make real-time decisions and coordinate actions efficiently. In this subsection, we delve into the mathematical principles and algorithmic strategies employed for adaptive scheduling in autonomous systems.

\subsubsection{Mathematically Detailed Discussion}
In autonomous systems, adaptive scheduling involves dynamically allocating tasks or resources based on changing environmental conditions, system states, and mission objectives. The scheduling problem often involves uncertainty and requires robust algorithms capable of adapting to unforeseen events and disturbances.

One adaptive scheduling approach commonly used in autonomous systems is \textbf{Reinforcement Learning (RL)}. RL formulates the scheduling problem as a Markov Decision Process (MDP) and learns optimal scheduling policies through trial and error interactions with the environment. The algorithm explores different scheduling strategies and updates its policy based on observed rewards and penalties.

\subsubsection{Algorithmic Example}
Here's a simple Python implementation of a Reinforcement Learning (RL) algorithm for adaptive scheduling in an autonomous system:
\FloatBarrier
\begin{lstlisting}
def reinforcement_learning():
    initialize_q_table()
    for episode in range(total_episodes):
        observe_environment()
        while not is_terminal_state():
            choose_action_based_on_policy()
            take_action()
            update_q_table()
\end{lstlisting}
\FloatBarrier

\section{Conclusion}
In this section, we summarize the key insights gained from exploring scheduling techniques related to algorithms and outline potential directions for future research and applications.

\subsection{The Evolution of Scheduling Theory and Practice}
Scheduling theory has undergone significant evolution over the years, driven by advances in computer science, optimization, and operations research. Initially, scheduling problems were primarily focused on job shop scheduling, where a set of jobs with specific processing requirements needed to be scheduled on available machines to optimize various objectives such as minimizing makespan or total completion time.

Over time, scheduling theory expanded to encompass a wide range of domains and problem settings, including project scheduling, task scheduling in parallel and distributed systems, resource allocation in cloud computing environments, and scheduling in manufacturing and transportation systems. This evolution was facilitated by the development of new algorithmic techniques, mathematical models, and optimization methods tailored to specific scheduling problems.

Early scheduling algorithms, such as the earliest versions of the List Scheduling algorithm for parallel machines, were often heuristic in nature, relying on simple rules or priority-based heuristics to make scheduling decisions. However, as computing power increased and optimization techniques advanced, more sophisticated algorithms based on mathematical programming, dynamic programming, and combinatorial optimization emerged, enabling the efficient solution of complex scheduling problems with large problem instances.

Today, scheduling theory and practice continue to evolve rapidly, driven by emerging technologies such as artificial intelligence, machine learning, and the Internet of Things. These technologies enable the development of adaptive and predictive scheduling algorithms that can dynamically adjust to changing conditions and optimize scheduling decisions in real-time. Additionally, interdisciplinary collaborations between researchers in computer science, operations research, and engineering are leading to the development of innovative scheduling solutions that address the unique challenges posed by modern applications and domains.

\subsection{Implications for Future Research and Applications}
The future of scheduling research holds immense promise, with numerous opportunities for advancing the state-of-the-art and addressing real-world challenges in various domains. Some implications for future research and applications include:

\begin{itemize}
    \item Development of intelligent scheduling algorithms: Future research will focus on developing intelligent scheduling algorithms that leverage techniques from artificial intelligence, machine learning, and data analytics to improve decision-making and optimize scheduling objectives.
    \item Integration of uncertainty and risk management: Scheduling algorithms will need to incorporate mechanisms for managing uncertainty and risk, enabling robust scheduling decisions in dynamic and uncertain environments.
    \item Scalability and efficiency: With the increasing scale and complexity of modern applications, there is a growing need for scalable and efficient scheduling algorithms that can handle large problem instances and optimize scheduling objectives in real-time.
\end{itemize}

In addition to research implications, scheduling algorithms have numerous applications across various domains, including:

\begin{itemize}
    \item Manufacturing: Scheduling algorithms are used to optimize production schedules, minimize setup times, and maximize resource utilization in manufacturing facilities.
    \item Transportation and logistics: Scheduling algorithms play a crucial role in optimizing routes, scheduling vehicle maintenance, and coordinating deliveries in transportation and logistics networks.
    \item Healthcare: Scheduling algorithms are used to optimize patient appointments, allocate resources, and manage healthcare facilities more efficiently.
\end{itemize}

These are just a few examples of the wide-ranging applications of scheduling algorithms, highlighting their importance in optimizing resource allocation and improving operational efficiency across diverse domains.

\section{Exercises and Problems}
In this section, we present a variety of exercises and problems related to scheduling algorithm techniques. These exercises aim to reinforce the concepts covered in the chapter and provide opportunities for students to apply their knowledge in practical scenarios. The section is divided into two subsections: Conceptual Questions to Test Understanding and Practical Scenarios for Algorithm Design.

\subsection{Conceptual Questions to Test Understanding}
This subsection contains a series of conceptual questions designed to test the reader's understanding of scheduling algorithm techniques. These questions cover key concepts, properties, and theoretical aspects of various scheduling algorithms.

\begin{itemize}
    \item Explain the difference between preemptive and non-preemptive scheduling algorithms.
    \item Describe the criteria used to evaluate the performance of scheduling algorithms.
    \item Discuss the advantages and disadvantages of First-Come, First-Served (FCFS) scheduling.
    \item Compare and contrast Shortest Job First (SJF) and Shortest Remaining Time (SRT) scheduling algorithms.
    \item Explain how Round Robin (RR) scheduling works and discuss its time quantum parameter.
\end{itemize}

\subsection{Practical Scenarios for Algorithm Design}
In this subsection, we present practical scenarios related to scheduling algorithm techniques. These scenarios involve real-world problems that require the design and implementation of scheduling algorithms to optimize resource allocation and improve system performance.

\begin{itemize}
    \item \textbf{Scenario 1: CPU Scheduling}
    
    You are tasked with implementing a CPU scheduling algorithm for a multi-tasking operating system. Design an algorithm that minimizes average waiting time and ensures fairness among processes.
    
    \begin{algorithm}[H]
        \caption{Shortest Job First (SJF) Scheduling}
        \begin{algorithmic}[1]
            \Procedure{SJF}{}
            \State Sort processes by burst time in ascending order
            \State Initialize total\_time = 0, total\_waiting\_time = 0
            \For{each process in sorted processes}
                \State total\_time += process.burst\_time
                \State total\_waiting\_time += total\_time - process.arrival\_time - process.burst\_time
            \EndFor
            \State average\_waiting\_time = total\_waiting\_time / number\_of\_processes
            \State \textbf{return} average\_waiting\_time
            \EndProcedure
        \end{algorithmic}
    \end{algorithm}
    
    \FloatBarrier
    \begin{lstlisting}[language=Python, caption=Python Implementation of SJF Scheduling Algorithm]
    def sjf_scheduling(processes):
        processes.sort(key=lambda x: x.burst_time)
        total_time = 0
        total_waiting_time = 0
        for process in processes:
            total_time += process.burst_time
            total_waiting_time += total_time - process.arrival_time - process.burst_time
        average_waiting_time = total_waiting_time / len(processes)
        return average_waiting_time
    \end{lstlisting}
    \FloatBarrier
    
    \item \textbf{Scenario 2: Real-Time Task Scheduling}
    
    Your task is to develop a real-time task scheduling algorithm for an embedded system. The system has multiple tasks with different execution times and deadlines. Design an algorithm that meets all task deadlines while maximizing system throughput.
    
    \begin{algorithm}[H]
        \caption{Rate-Monotonic Scheduling}
        \begin{algorithmic}[1]
            \Procedure{RateMonotonic}{}
            \State Sort tasks by their periods in ascending order
            \State Assign priorities to tasks based on their periods (shorter periods have higher priorities)
            \State Execute tasks in order of priority, preempting lower-priority tasks if necessary
            \EndProcedure
        \end{algorithmic}
    \end{algorithm}
    
    \FloatBarrier
    \begin{lstlisting}[language=Python, caption=Python Implementation of Rate-Monotonic Scheduling Algorithm]
    def rate_monotonic_scheduling(tasks):
        tasks.sort(key=lambda x: x.period)
        for task in tasks:
            execute_task(task)
    \end{lstlisting}
    \FloatBarrier
    
    \item \textbf{Scenario 3: Disk Scheduling}
    
    Develop a disk scheduling algorithm for a hard disk drive to optimize disk access time. The algorithm should minimize disk head movement and ensure fair access to disk requests.
    
    \begin{algorithm}[H]
        \caption{SCAN (Elevator) Disk Scheduling}
        \begin{algorithmic}[1]
            \Procedure{SCAN}{}
            \State Sort disk requests by cylinder number
            \State Initialize head direction to right
            \State Traverse requests in sorted order, servicing requests until the end
            \State Change head direction when reaching the last request
            \State Continue traversing in the opposite direction until all requests are serviced
            \EndProcedure
        \end{algorithmic}
    \end{algorithm}
    
    \FloatBarrier
    \begin{lstlisting}[language=Python, caption=Python Implementation of SCAN Disk Scheduling Algorithm]
    def scan_disk_scheduling(requests):
        requests.sort()
        head_direction = 'right'
        while requests:
            if head_direction == 'right':
                for request in requests:
                    service_request(request)
                head_direction = 'left'
            else:
                for request in reversed(requests):
                    service_request(request)
                head_direction = 'right'
        \end{lstlisting}
    \FloatBarrier
\end{itemize}

\section{Further Reading and Resources}

When diving deeper into scheduling algorithm techniques, it's essential to explore various resources to gain a comprehensive understanding. Below are some valuable resources categorized into foundational texts and influential papers, online courses and tutorials, and software and tools for scheduling analysis.

\subsection{Foundational Texts and Influential Papers}

Understanding the foundational principles and key research papers is crucial for mastering scheduling algorithm techniques. Here's a curated list of influential texts and papers:

\begin{itemize}
    \item \textbf{Foundational Texts}:
    \begin{itemize}
        \item "Introduction to Algorithms" by Thomas H. Cormen, Charles E. Leiserson, Ronald L. Rivest, and Clifford Stein.
        \item "Scheduling: Theory, Algorithms, and Systems" by Michael Pinedo.
        \item "Operating System Concepts" by Abraham Silberschatz, Peter Baer Galvin, and Greg Gagne.
    \end{itemize}
    
    \item \textbf{Influential Papers}:
    \begin{itemize}
        \item "Optimal Control of Discrete-Event Systems" by N. N. Ulanov, M. I. Kudenko, and G. V. Tsarev.
        \item "Analysis of Multilevel Feedback Queues" by W. A. Halang and G. B. Halang.
        \item "Real-Time Scheduling Theory: A Historical Perspective" by C. L. Liu and J. W. Layland.
    \end{itemize}
\end{itemize}

\subsection{Online Courses and Tutorials}

Online courses and tutorials offer interactive learning experiences and in-depth explanations of scheduling algorithms. Here are some recommended resources:

\begin{itemize}
    \item \textbf{Coursera}:
    \begin{itemize}
        \item "Operating Systems" by University of London.
        \item "Discrete Optimization" by The University of Melbourne.
    \end{itemize}
    
    \item \textbf{edX}:
    \begin{itemize}
        \item "Algorithm Design and Analysis" by University of Pennsylvania.
        \item "Real-Time Systems" by Karlsruhe Institute of Technology.
    \end{itemize}
    
    \item \textbf{YouTube Tutorials}:
    \begin{itemize}
        \item "Introduction to Scheduling Algorithms" by MyCodeSchool.
        \item "Real-Time Operating Systems Concepts" by Gate Lectures by Ravindrababu Ravula.
    \end{itemize}
\end{itemize}

\subsection{Software and Tools for Scheduling Analysis}

Utilizing specialized software and tools can streamline scheduling analysis tasks and aid in implementing and testing scheduling algorithms. Here are some popular options:

\begin{itemize}
    \item \textbf{SimPy}: SimPy is a discrete-event simulation library for Python. It provides components for modeling and simulating various scheduling scenarios.
    
    \item \textbf{GanttProject}: GanttProject is an open-source project scheduling and management tool. It allows users to create Gantt charts to visualize and manage scheduling tasks.
    
    \item \textbf{MATLAB}: MATLAB offers robust tools for numerical computing and simulation. It includes functions and toolboxes for modeling and analyzing scheduling algorithms.
    
    \item \textbf{SimEvents}: SimEvents is a discrete-event simulation toolset in MATLAB that extends Simulink for modeling and analyzing event-driven systems, including scheduling processes.
\end{itemize}
\end{document}